\section{Everett Interpretation}

The mathematical framework becomes an Everettian observer-selection model only after an additional physical identification. We state this explicitly as the Everett-QBS bridge assumption:
\[
d\mu^{\mathrm{FP}}_\pi(\omega)
=
\frac{S_\pi(\omega)}{\mathbb E_\mu[S_\pi]}
\,d\mu(\omega).
\]
Under this assumption, policy-dependent accessibility acts as an observer-indexed self-location weight over the base branch measure.

The present work does not derive this bridge from unitary quantum mechanics, decoherence theory, or the Born rule. The bridge may therefore be accepted, modified, or rejected independently of the measure-theoretic theorems. Rejecting it removes the Everettian physical interpretation but leaves the probability identities and classical policy simulations intact.

This separation is methodologically important. The core claims concern how a weighted first-person measure behaves once an accessibility function is specified and how recognition-dependent policies can change both trajectories and accessibility. The unresolved physical claim is whether Everettian first-person uncertainty is represented by such a policy-dependent accessibility function.

The manuscript will position this bridge relative to prior work on Everettian probability, self-locating uncertainty, and anthropic/observer-selection decision theory. That literature review remains in progress during the public-review phase.
